Durante o desenvolvimento de aplicações mais complexas é comum a utilização 
de bibliotecas para facilitar o processo de desenvolvimento. Portanto, iremos 
cumprir aqui neste capítulo o básico que você precisa saber para lidar com a 
utilização e implementação de bibliotecas em um projeto.
\\\\
Para entender melhor sobre o assunto vamos usar o seguinte código como exemplo:

\begin{verbatim}
    #include <stdio.h>
    #include <math.h>

    int main(int argc, char **argv)
    {
        // Armazena os valores
        int c1 = 3;
        int c2 = 4;

        // Realiza os cálculos
        float h = sqrt(pow(c1, 2) + pow(c2, 2));

        // Exibir resultados
        printf("Valor da hipotenusa: %.2f\n", h);
        return 0;
    }
\end{verbatim}

O código acima é responsável por calcular a hipotenusa de um triângulo 
conforme os valores dados pelas variáveis \(c1\) e \(c2\). Porém, os mais 
atentos notaram o uso das funções \textbf{pow} e \textbf{sqrt}, responsáveis 
por calcular a potência e a raiz quadrada de um número, respectivamente.
\\\\
Ambas as funções estão inclusas na biblioteca \textbf{math.h} que não é uma 
biblioteca padrão da linguagem C. Dito isso, caso você tente digitar compilar 
o código acima usando o código:

\begin{verbatim}
    gcc -o main main.c
\end{verbatim}

É provável que você veja um erro semelhante ao seguinte:

\begin{verbatim}
    /usr/bin/ld: /tmp/ccF2H2AC.o: in function `main':
    main.c:(.text+0x3d): undefined reference to `pow'
    /usr/bin/ld: main.c:(.text+0x66): undefined reference to `pow'
    /usr/bin/ld: main.c:(.text+0x7e): undefined reference to `sqrt'
    collect2: error: ld returned 1 exit status
\end{verbatim}

Isso indica que, embora a biblioteca tenha sido incluída adequadamente no 
código, ela não está sendo linkada pelo compilador, resultando na não 
localização das funções.
\\\\
Para corrigir isso, basta informar ao compilador que você pretende compilar 
o código usando uma biblioteca específica. Neste caso, como estamos incluindo 
a biblioteca \texttt{math}, devemos adicionar ``-lm'' ao comando de compilação, 
como mostrado abaixo:

\begin{verbatim}
    gcc -o main main.c -lm
\end{verbatim}

Dessa forma, o compilador reconhecerá a biblioteca e realizará a compilação 
corretamente.

\subsubsection{Exercícios}

\textbf{É de extrema importância que você realize os exercícios aqui elaborados para fixar o 
conteúdo.}

\begin{enumerate}
    \item{Tente compilar o código acima sem valor de `y'.}
\end{enumerate}

\clearpage
